\cleardoublepage

\section{软件开发文档}

% ============================================================
% 本章约4500-6000字，对应毕业设计要求的"软件开发文档"
% 作为第4章
% 内容包括：软件设计文档、使用说明书、测试分析报告等
% ============================================================

\subsection{API接口设计文档}

% -----------------------------------------------------------
% 内容要点：
%   - 介绍系统的RESTful API设计规范和接口列表
%   - 可引用Swagger生成的API文档
% -----------------------------------------------------------

\subsubsection{接口设计规范}

\par 本项目的对外接口遵循 RESTful 风格，以资源为核心，以 HTTP 方法表达动作。所有接口路径统一以 \texttt{/api} 为前缀，下文为简洁起见均省略。该前缀通过框架路径匹配配置器动态拼接，无需修改控制器代码即可调整对外暴露的根路径。

\par 资源命名采用单复数二元约定：操作单个资源使用单数名词加路径参数，例如 GET、PATCH、DELETE \texttt{/project/\{id\}}；获取列表使用复数名词，例如 GET \texttt{/projects}。多单词资源名采用短横线连接，如 \texttt{/project-view}。存在父子层级的资源在路径中显式嵌入父资源标识符后再追加子资源段，如 \texttt{/project/\{id\}/task} 与 \texttt{/session/\{id\}/chats}。

\par 系统严格遵循 HTTP 方法语义：GET 用于无副作用的查询，参数通过查询字符串传递；POST 用于创建资源或触发不可幂等动作，请求体以 JSON 承载；PATCH 用于局部更新，仅传入需修改字段；DELETE 用于删除资源。考虑到 AI 智能体常需一次性写入多条任务，创建、更新和删除接口均同时提供单条与批量两种变体，批量路径追加 \texttt{/batch} 后缀，由事务保证全部成功或全部回滚。

\par 请求与响应采用 DTO 与 VO 分离模式。写操作接口以专门的 DTO 承载入参并通过声明式校验注解进行前置约束；成功响应以专门的 VO 类型承载，避免将持久化对象直接暴露给前端。所有响应统一包裹为 \texttt{Response<T>}，包含状态码、提示消息和数据载荷。状态码采用五位数字编码，前两位表示错误类别，后三位表示具体错误，如 \texttt{20001} 表示成功，\texttt{40101} 表示未授权，每个错误码关联到标准 HTTP 状态码并由全局异常处理器同步设置到响应头。

\par 系统采用基于 JWT 的无状态认证方案。客户端登录后获得访问令牌，后续请求以 \texttt{Bearer <token>} 形式置于 \texttt{Authorization} 头中传递。安全过滤器在请求进入业务逻辑前解析验证令牌，包括签名合法性、有效期与黑名单状态；校验通过后将用户身份绑定到认证上下文，否则返回 \texttt{40101} 与 \texttt{401}。登录与验证码接口以及运维探活端点配置为白名单放行。

\par 跨域访问通过 CORS 统一配置，对所有路径放行 GET、POST、PATCH、DELETE 和 OPTIONS 方法。AI 对话接口采用 Server-Sent Events 流式响应，按模型生成节奏向前端推送增量文本与工具调用结果，客户端持续读取并增量渲染，最后一个事件标记为 \texttt{done}。

\par 接口契约通过 SpringDoc OpenAPI 自动生成，以 Swagger UI 形式暴露在 \texttt{/swagger-ui}，文档包含接口路径、请求响应模型、字段校验约束与示例值。规范文件以 JSON 发布于 \texttt{/v3/api-docs}，便于代码生成工具消费。

\subsubsection{核心接口列表}

\par 系统对外暴露的接口超过六十个，按业务域划分为用户认证、项目管理、智能体对话、语音服务与运维监控五个子域。下文按子域分别列出核心接口，每个接口给出 HTTP 方法、路径与简要说明，完整的请求体与响应体定义可在系统部署后访问 \texttt{/swagger-ui} 路径查看在线 OpenAPI 文档。

\par 用户认证模块承担账号创建、身份核验与会话失效三类职责，由于其中部分接口必须在未认证状态下可被调用，安全模块在过滤器链中为登录与发送验证码两个端点配置了显式的白名单放行规则。表 \ref{tab:api-user} 列出该模块的全部接口。

\begin{table}[htbp]
  \centering
  \caption{用户认证模块接口列表}
  \label{tab:api-user}
  \begin{tabular}{p{2cm}p{5cm}p{6cm}}
    \hline
    方法 & 路径 & 说明 \\
    \hline
    POST & \texttt{/user/send-code} & 向指定手机下发六位短信验证码 \\
    POST & \texttt{/user/login} & 凭手验证码登录并注册 \\
    POST & \texttt{/user/logout} & 将当前访问令牌加入登出黑名单 \\
    GET & \texttt{/user/\{id\}} & 查询指定用户的资料 \\
    PATCH & \texttt{/user/\{id\}} & 修改昵称、头像等个人资料 \\
    DELETE & \texttt{/user/\{id\}} & 注销账号并清理其名下数据 \\
    \hline
  \end{tabular}
\end{table}

\par 项目与任务接口构成任务管理域的基础数据结构，是用户规划与追踪进度的核心载体。每种资源均同时提供单条与批量两种变体，批量变体的路径在原路径基础上追加 \texttt{/batch} 后缀，请求体携带数组形式的操作载荷，由后端在事务内统一执行，便于 AI 智能体在一次任务拆解中一次性写入多条任务。任务记录标题、截止日期、优先级与完成状态等字段，用户可通过状态筛选快速定位当前关注的工作项。任务隶属于具体项目，因此创建任务和按项目分页查询任务的路径以项目标识符为父路径段，而对任务自身的更新和删除则使用扁平路径 \texttt{/task/\{id\}}。表 \ref{tab:api-project} 汇总了该模块的接口。

\begin{table}[htbp]
  \centering
  \caption{项目与任务管理模块接口列表}
  \label{tab:api-project}
  \begin{tabular}{p{2cm}p{5cm}p{6cm}}
    \hline
    方法 & 路径 & 说明 \\
    \hline
    GET & \texttt{/projects} & 分页查询当前用户的项目列表 \\
    GET & \texttt{/project/\{id\}} & 查询单个项目的详情 \\
    POST & \texttt{/project} & 创建单个项目 \\
    POST & \texttt{/project/batch} & 批量创建项目 \\
    PATCH & \texttt{/project/\{id\}} & 局部更新单个项目 \\
    PATCH & \texttt{/project/batch} & 批量更新项目 \\
    DELETE & \texttt{/project/\{id\}} & 删除单个项目并级联清理 \\
    DELETE & \texttt{/project/batch} & 批量删除项目 \\
    GET & \texttt{/project/\{id\}/tasks} & 分页查询指定项目下的任务 \\
    POST & \texttt{/project/\{id\}/task} & 在指定项目下创建任务 \\
    POST & \texttt{/project/\{id\}/task/batch} & 在指定项目下批量创建任务 \\
    GET & \texttt{/task/\{id\}} & 查询单个任务的详情 \\
    PATCH & \texttt{/task/\{id\}} & 局部更新单个任务 \\
    PATCH & \texttt{/task/batch} & 批量更新任务 \\
    DELETE & \texttt{/task/\{id\}} & 删除单个任务 \\
    DELETE & \texttt{/task/batch} & 批量删除任务 \\
    \hline
  \end{tabular}
\end{table}

\par 标签与项目—标签关联模块为项目的二次组织提供能力，便于用户按自定义维度对项目进行分类与检索。标签自身是独立的资源，拥有独立的标识符、名称与颜色字段，用户可为紧急事项创建红色高优标签，或为长期目标创建蓝色规划标签，通过颜色快速识别项目属性。标签遵循与项目相同的增删改查接口约定；项目与标签之间的多对多关联则通过专门的 \texttt{/project-tag} 资源管理，支持以项目为视角查询其全部标签，或以标签为视角查询其下挂载的全部项目。将关联关系独立为资源而非在项目表中嵌入标签列表，既保留了标签的复用性，也便于前端以统一方式处理关联的增删改查与批量操作，同时避免了项目表因标签字段膨胀而导致的结构不稳定。表 \ref{tab:api-tag} 列出该模块的接口。

\begin{table}[htbp]
  \centering
  \caption{标签与项目-标签关联模块接口列表}
  \label{tab:api-tag}
  \begin{tabular}{p{2cm}p{5cm}p{6cm}}
    \hline
    方法 & 路径 & 说明 \\
    \hline
    GET & \texttt{/tags} & 分页查询当前用户的标签列表 \\
    GET & \texttt{/tag/\{id\}} & 查询单个标签的详情 \\
    POST & \texttt{/tag} & 创建单个标签 \\
    POST & \texttt{/tag/batch} & 批量创建标签 \\
    PATCH & \texttt{/tag/\{id\}} & 局部更新单个标签 \\
    PATCH & \texttt{/tag/batch} & 批量更新标签 \\
    DELETE & \texttt{/tag/\{id\}} & 删除单个标签 \\
    DELETE & \texttt{/tag/batch} & 批量删除标签 \\
    GET & \texttt{/project-tag} & 查询单条项目—标签关联信息 \\
    POST & \texttt{/project-tag} & 创建单条项目—标签关联 \\
    POST & \texttt{/project-tag/batch} & 批量创建项目—标签关联 \\
    DELETE & \texttt{/project-tag} & 删除单条项目—标签关联 \\
    DELETE & \texttt{/project-tag/batch} & 批量删除项目—标签关联 \\
    GET & \texttt{/project/\{id\}/tags} & 查询指定项目挂载的全部标签 \\
    GET & \texttt{/tag/\{id\}/projects} & 查询指定标签下的全部项目 \\
    \hline
  \end{tabular}
\end{table}

\par 项目视图是面向前端首页一站式渲染以及 AI 智能体上下文输入的聚合查询资源。每条视图记录捕获项目主信息与任务列表、标签列表的快照，前端可以一次性获取展示所需的全部字段，而无需多次往返调用项目、任务和标签三个独立接口。由于该资源同时承担读写两种语义——既允许前端订阅服务端聚合后的视图，也允许 AI 智能体在工具调用中直接同步草稿到视图字段——更新接口在语义上命名为同步而非更新。表 \ref{tab:api-projectview} 列出该模块的接口。

\begin{table}[htbp]
  \centering
  \caption{项目视图模块接口列表}
  \label{tab:api-projectview}
  \begin{tabular}{p{2cm}p{5cm}p{6cm}}
    \hline
    方法 & 路径 & 说明 \\
    \hline
    GET & \texttt{/project-views} & 分页查询当前用户的项目视图 \\
    GET & \texttt{/project-view/\{id\}} & 查询单个项目视图的详情 \\
    POST & \texttt{/project-view} & 创建单个项目视图记录 \\
    POST & \texttt{/project-view/batch} & 批量创建项目视图记录 \\
    PATCH & \texttt{/project-view/\{id\}} & 同步单个项目视图的草稿字段 \\
    PATCH & \texttt{/project-view/batch} & 批量同步项目视图草稿字段 \\
    DELETE & \texttt{/project-view/\{id\}} & 删除单个项目视图记录 \\
    DELETE & \texttt{/project-view/batch} & 批量删除项目视图记录 \\
    \hline
  \end{tabular}
\end{table}

\par 会话与消息接口承载智能体对话的元数据持久化。一个会话对应一段对话上下文，可以绑定到一个具体项目，也可以独立存在；绑定项目后，智能体在回复时可以直接引用该项目的任务列表与标签信息作为上下文，从而给出更具针对性的建议，同一项目下的多轮对话由此保持上下文连续性。会话下的消息按时间顺序排列，每条消息记录角色标识与内容文本，角色分为用户与智能体两类，构成完整的消息列表。前端在打开对话窗口时通过会话接口加载或创建会话，再通过消息接口分页拉取历史消息进行渲染，避免一次性加载过多消息导致前端性能下降。会话的创建时间与最后活跃时间由服务端自动维护，便于前端按最近使用排序展示。消息隶属于会话，因此创建消息和按会话分页拉取消息的路径以会话标识符为父路径段，而消息自身的更新和删除则使用扁平路径 \texttt{/chat/\{id\}}。

\par 智能体模块依赖会话与消息接口协同工作：会话接口提供对话上下文，消息接口记录交互历史，智能体接口读取历史消息作为模型输入。三者通过会话标识符串联，形成对话闭环。

\par 表 \ref{tab:api-session} 列出该模块的接口。

\begin{table}[htbp]
  \centering
  \caption{会话与消息管理模块接口列表}
  \label{tab:api-session}
  \begin{tabular}{p{2cm}p{5cm}p{6cm}}
    \hline
    方法 & 路径 & 说明 \\
    \hline
    GET & \texttt{/sessions} & 分页查询当前用户的会话列表 \\
    GET & \texttt{/session/\{id\}} & 查询单个会话的元数据 \\
    POST & \texttt{/session} & 创建单个会话 \\
    POST & \texttt{/session/batch} & 批量创建会话 \\
    PATCH & \texttt{/session/\{id\}} & 局部更新单个会话 \\
    PATCH & \texttt{/session/batch} & 批量更新会话 \\
    DELETE & \texttt{/session/\{id\}} & 删除单个会话 \\
    DELETE & \texttt{/session/batch} & 批量删除会话 \\
    GET & \texttt{/session/\{id\}/chats} & 分页查询指定会话下的消息 \\
    POST & \texttt{/session/\{id\}/chat} & 在指定会话下追加一条消息 \\
    POST & \texttt{/session/\{id\}/chat/batch} & 在指定会话下批量追加消息 \\
    GET & \texttt{/chat/\{id\}} & 查询单条消息的详情 \\
    PATCH & \texttt{/chat/\{id\}} & 局部更新单条消息 \\
    PATCH & \texttt{/chat/batch} & 批量更新消息 \\
    DELETE & \texttt{/chat/\{id\}} & 删除单条消息 \\
    DELETE & \texttt{/chat/batch} & 批量删除消息 \\
    \hline
  \end{tabular}
\end{table}

\par 智能体对话接口是整个系统中唯一的流式接口，响应体不再是单一的 JSON 对象，而是遵循 Server-Sent Events 协议的事件流，按模型生成节奏向前端推送增量文本与工具调用结果。该接口提供两种入口：在已有会话上下文中追加用户输入并继续对话的有状态入口 \texttt{/session/\{id\}/agent/chat}，以及无需会话标识符即可发起独立对话的无状态入口 \texttt{/agent/chat}。当模型在生成过程中触发工具调用时，服务端会暂停输出并执行相应操作，待获取结果后再将上下文重新注入模型继续生成，整个工具链的调用过程对前端透明。对话状态通过 Redis 持久化，支持用户在中断后随时追问而无需重新建立上下文。

\par 语音模块提供语音识别与语音合成两项能力，前端通过调用临时令牌签发接口获取访问凭证后，可直接连接阿里云语音服务进行实时语音转文字或文字转语音操作。令牌采用短期有效策略，过期后前端需重新申请，服务端在签发时根据配置动态拼接访问密钥与过期时间戳，确保即使令牌被截获也不会造成长期安全隐患。这一设计将长期访问凭证完全隔离在后端，避免密钥以任何形式暴露在客户端代码或网络传输中。考虑到临时令牌在有效期内可重复使用，服务端将已签发的令牌以语音服务凭证为维度缓存到 Redis，避免并发请求触发多次远程签发。后续请求命中缓存时直接返回已有令牌，显著降低令牌签发接口的平均响应延迟，同时减少对阿里云 NLS 服务的调用频次。

\par 运维监控接口由 Spring Boot Actuator 自动暴露，用于负载均衡器健康探活与发布版本核对。健康检查端点返回服务本身及其依赖组件的运行状态，信息端点则返回应用的构建版本号、Git 提交哈希与启动时间戳，便于运维人员在多实例部署环境中快速核对线上运行的代码版本。

\par 表 \ref{tab:api-agent} 列出这一组接口。

\begin{table}[htbp]
  \centering
  \caption{智能体对话与基础服务接口列表}
  \label{tab:api-agent}
  \begin{tabular}{p{2cm}p{5cm}p{6cm}}
    \hline
    方法 & 路径 & 说明 \\
    \hline
    POST & \texttt{/agent/chat} & 发起一次智能体对话，流式返回 \\
    POST & \texttt{/session/\{id\}/agent/chat} & 在指定会话下进行追问 \\
    POST & \texttt{/voice/token} & 签发用于语音服务的临时令牌 \\
    GET & \texttt{/actuator} & 列出全部已暴露的运维端点 \\
    GET & \texttt{/actuator/health} & 探查服务及其依赖的健康状态 \\
    GET & \texttt{/actuator/info} & 返回构建版本与启动信息 \\
    \hline
  \end{tabular}
\end{table}

\subsection{系统使用说明书}

\subsubsection{运行环境要求}

\par 本项目采用容器化部署，需根据用户规模预留处理器、内存和存储资源。系统依赖大语言模型接口，网络带宽应保证稳定，避免因网络抖动导致智能体对话出现明显卡顿。生产环境通过容器编排平台部署数据库、缓存与后端服务；开发环境额外需要 Java 开发工具包与构建工具。容器镜像通过多阶段构建控制体积，操作系统以 Linux 发行版为佳。存储方面，由于数据库容器使用绑定挂载持久化数据，建议为数据目录配置高速存储介质，避免机械硬盘读写延迟成为瓶颈。

\par 详细环境要求见附录。

\subsubsection{环境配置}

\par 系统采用环境变量文件驱动配置，所有敏感信息与运行时参数均通过根目录下的环境变量文件注入。环境变量按功能域分为数据库、缓存、安全、人工智能、短信和语音六大类。生产环境中环境变量文件应设置严格的文件权限，避免其他用户读取敏感信息。配置加载采用分层覆盖机制，系统启动时会先加载内置的默认配置，随后通过环境变量进行覆盖，优先级由低到高依次为默认值、配置文件、环境变量，保证配置在不同环境下的灵活切换。

\par 具体环境变量见附录。

\subsubsection{部署步骤}

\par 生产环境部署采用容器化一键部署方案。部署前完成仓库访问认证，随后执行远程脚本自动完成环境创建、镜像拉取与服务启动。部署完成后通过运维监控端点验证服务状态。

\par 本地开发部署同样简便。编译源码后启动基础设施容器，再启动后端应用即可。开发环境支持热部署，代码修改后无需手动重启。

\par 具体的生产部署或本地开发命令见附录。

\subsection{功能测试}

\subsubsection{测试环境}

\par 功能测试基于 Spring Boot 测试框架与 JUnit 5构建，测试数据库和缓存分别使用本地 PostgreSQL 和 Redis 实例。测试过程中通过清理工具自动清除数据库与缓存中的残留数据，保证测试之间的独立性。

\subsubsection{测试方法}

\par 普通接口测试采用接口级别的端到端测试方法。测试流程遵循“创建-查询-同步-删除”的完整生命周期，覆盖单条和批量两种操作模式。每个测试用例通过模拟 HTTP 请求验证响应状态码和返回数据的一致性，并通过递归比较忽略动态生成字段（如标识符和时间戳）来断言业务数据的正确性。批量操作测试验证事务一致性，确保全部成功或全部回滚。边界测试则验证查询不存在资源时返回404等异常情况。

\par 智能体测试采用流式响应订阅方式验证对话功能。测试覆盖智能路由、历史记忆加载和新建对话等场景。通过订阅服务端推送的流式事件，验证会话标识符的一致性、对话状态的持久化以及工具调用结果的正确性。项目构建测试分别针对简单和复杂两种项目类型，验证智能体能否正确解析用户意图、生成任务列表并完成标签关联，同时验证会话在多次对话间的上下文连续性。

\par 云服务集成测试采用真实模式调用云服务。测试覆盖成功请求、成功调用和错误调用等场景。通过配置测试环境连接真实第三方接口，验证系统在真实网络环境下的连通性、数据序列化和错误处理能力，确保与外部云服务的集成稳定可靠。

\subsubsection{测试结果}

\par 普通接口测试验证了各业务模块的完整生命周期与事务一致性。项目与任务模块测试覆盖单条与批量两种操作模式，批量操作在事务内统一执行，通过对比操作前后的列表查询结果验证事务一致性；边界测试验证查询不存在资源时返回404、越权操作返回403，列表分页返回结果与预期完全一致，未出现资源泄漏或状态异常。用户模块测试验证了登录鉴权、令牌失效与注销机制，错误令牌、越权请求和已注销令牌均被正确拦截。仓储集成测试对全部核心实体的单条与批量增删改查进行了完整覆盖，验证了各层级关联关系的正确性。

\par 智能体测试采用流式响应订阅方式，覆盖了智能路由、历史记忆和新建对话三个核心场景。测试验证了会话标识符在多轮交互间保持一致，对话状态通过Redis正确持久化；针对简单和复杂两种项目类型，分别验证了智能体能够根据用户意图生成相应任务列表并完成标签关联——简单项目生成微步骤，复杂项目生成三段式拆解，工具调用结果与预期一致。

\par 云服务集成测试在真实网络环境下验证了与第三方服务的连通性与错误处理能力。短信服务测试验证了真实验证码发送、登录流程和错误验证码被拒绝的场景；语音服务测试验证了临时令牌获取、Redis缓存命中与缓存失效后重新生成的机制，所有模块整体通过率达到100\%。

\begin{table}[htbp]
  \centering
  \caption{功能测试结果汇总}
  \label{tab:func-test-result}
  \begin{tabular}{p{4cm}p{4cm}}
    \hline
    指标 & 结果 \\
    \hline
    测试用例总数 & 123 \\
    整体通过率 & 100\% \\
    \hline
  \end{tabular}
\end{table}

\subsection{性能测试}

\subsubsection{测试环境}

\par 性能测试基于 Spring Boot 测试框架构建，运行环境为 2 核 4 线程 16 GB 内存，数据库和缓存服务均通过 Docker 容器化运行（PostgreSQL 15、Redis 7），保证测试环境与生产环境一致。

\subsubsection{测试方法}

\par 接口响应维度针对创建项目等核心端点执行多次请求，用高精度计时器采样 P50、P95、P99 延迟并统计错误率；同时通过多线程在固定时间窗口内持续发起请求，统计每秒请求数作为实际吞吐量；最后基于线程池规模和平均延迟进行理论最大吞吐量估算，为容量规划提供参考。

\par 数据库与 JVM 维度包含多项细粒度指标。数据库方面，通过对比批量创建项目前后列表查询的延迟增长幅度来定性检测是否存在 N+1 查询问题；同时独立执行 Flyway 迁移脚本并统计总耗时，以评估数据库初始化效率。JVM 方面，在批量创建项目前后采样堆内存峰值，并在过程中定期采样绘制增长曲线，观察内存使用的线性趋势；通过对比公开端点与受保护端点的延迟差值评估 JWT 认证开销；对提示词模板执行多次字符串替换操作，评估微秒级的渲染耗时。

\par 并发稳定性维度通过多线程同步启动同时修改同一项目，验证数据一致性与是否存在死锁；并在并发测试过程中定期采样线程数量，观察线程池是否出现饱和，以评估系统在并发压力下的资源利用情况。

\par 智能体维度针对流式对话场景进行专项测试。通过模拟延迟流测量从发起请求到首字节返回的时间；模拟固定间隔的数据块下发，统计流式输出的平均间隔时间；模拟异常返回场景，测量降级快速失败的延迟；调用对话状态持久化接口，测量 Redis 写入延迟；模拟固定长度的响应内容，结合估算的 token 数量计算 token 吞吐量。

\par 缓存维度针对语音服务令牌缓存进行测试。通过先后触发缓存未命中和缓存命中两种场景，对比同一接口在两种状态下的延迟差异，从而计算缓存加速比，验证缓存对高频访问接口的实际加速效果。

\subsubsection{测试结果}

\par 从接口响应维度来看，核心端点的 P50 与 P95 延迟均达到了毫秒级水平，能够满足日常交互中对响应速度的苛刻要求；理论最大吞吐量达到了万级 RPS 水平，表明系统在并发承载方面具备较强的潜力。

\par 从数据库与 JVM 维度来看，列表查询的延迟增长比例处于较低水平，说明在当前数据规模下不存在显著的 N+1 查询问题，ORM 的抓取策略与索引设计基本合理；Flyway 迁移耗时达到了亚秒级水平，能够在应用启动阶段快速完成 schema 初始化，满足持续交付场景下对部署效率的要求；JWT 认证与数据库访问的综合开销达到了可控水平，说明安全校验与数据持久化环节并未成为整体性能的瓶颈。

\par 从并发稳定性维度来看，多线程并发更新全部成功且无异常，数据一致性和事务隔离达到了预期水平，表明悲观锁或乐观锁的配置能够有效避免竞态条件；线程池未出现饱和迹象，说明当前并发压力尚未触及资源瓶颈，系统仍具备一定的弹性扩展空间。后续随着用户规模增长，建议持续监控线程池队列长度与活跃线程数，以便及时调整线程池参数。

\par 从智能体维度来看，首字节时间（TTFB）达到了可接受的水平，用户发起对话后能够在较短时间内收到首条响应，基本满足实时交互的体验要求；流式输出间隔达到了稳定水平，说明服务端事件推送节奏均匀，前端渲染不会出现明显卡顿；降级快速失败机制的响应速度达到了预期水平，能够在模型服务异常时迅速切换至兜底逻辑，避免用户长时间等待；对话状态持久化的延迟达到了毫秒级水平，Redis 写入性能良好，为断线重连与历史会话恢复提供了可靠保障。

\par 从缓存维度来看，Redis 缓存对语音服务令牌的加速效果达到了显著水平，缓存命中后的接口延迟从数百毫秒级降至毫秒级，加速比达到了数十倍水平，充分说明缓存策略对高频访问接口的提升效果十分明显。考虑到语音令牌接口在对话高峰期可能被集中调用，建议进一步监控缓存命中率与过期时间设置，以平衡实时性与系统负载。

\par 需要说明的是，所有延迟指标均为本地测试环境所得，与真实生产网络存在差异，实际部署时应结合负载均衡与横向扩容策略进一步验证；此外，Token 吞吐量指标反映的是框架层输出速度，而非模型推理速度，因此在评估端到端对话体验时，仍需将模型自身的生成延迟纳入考量。

\begin{table}[htbp]
  \centering
  \caption{性能测试结果汇总}
  \label{tab:perf-summary}
  \begin{tabular}{p{4cm}p{4cm}}
    \hline
    指标 & 结果 \\
    \hline
    \multicolumn{2}{l}{\textbf{接口响应}} \\
    P50 延迟 & 7.991 ms \\
    P95 延迟 & 11.129 ms \\
    P99 延迟 & 13.712 ms \\
    平均延迟 & 8.160 ms \\
    理论最大吞吐量 & 25,026.9 RPS \\
    \hline
    \multicolumn{2}{l}{\textbf{数据库与 JVM}} \\
    N+1 查询检测 & 延迟比例 0.45x \\
    Flyway 迁移耗时 & 156 ms \\
    JWT 与 DB 开销 & 7.576 ms \\
    \hline
    \multicolumn{2}{l}{\textbf{并发稳定性}} \\
    并发更新线程数 & 8 \\
    成功数/异常数 & 8/0 \\
    线程数范围 & 35--45 \\
    采样点数 & 51 \\
    \hline
    \multicolumn{2}{l}{\textbf{智能体}} \\
    首字节时间（TTFB） & 61.113 ms \\
    流式输出间隔 & 18.505 ms \\
    降级快速失败延迟 & 20.040 ms \\
    对话状态持久化延迟 & 56.093 ms \\
    Token 吞吐量 & 264.6 ktps \\
    \hline
    \multicolumn{2}{l}{\textbf{缓存}} \\
    缓存未命中延迟 & 357.111 ms \\
    缓存命中延迟 & 4.969 ms \\
    缓存加速比 & 71.87x \\
    \hline
  \end{tabular}
\end{table}

\subsection{CI/CD流水线}

\subsubsection{测试环境}

\par 持续集成与持续部署流水线基于 GitHub Actions 构建，运行环境为最新版 Ubuntu。测试所需的敏感环境变量通过 GitHub 密钥库注入。

\subsubsection{测试方法}

\par 流水线在每次拉取请求时自动触发，功能测试全部通过是合并的必要条件。执行流程包括启动 PostgreSQL 和 Redis 服务容器、初始化数据库、执行测试、生成覆盖率报告、构建 Docker 镜像并推送到容器仓库。容器镜像采用多阶段构建策略，构建阶段基于 JDK 21完成编译打包，运行阶段基于 JRE 21精简镜像，创建非 root 用户运行应用，提升安全性。

\subsubsection{测试结果}

\par 任何测试失败都会阻止镜像构建与部署流程，确保只有经过充分验证的代码才能进入生产环境。

\begin{table}[htbp]
  \centering
  \caption{CI/CD 流水线执行结果}
  \label{tab:cicd-result}
  \begin{tabular}{p{4cm}p{4cm}}
    \hline
    指标 & 结果 \\
    \hline
    触发条件 & 每次拉取请求 \\
    功能测试通过 & 合并必要条件 \\
    覆盖率报告生成 & JaCoCo \\
    镜像构建策略 & 多阶段构建 \\
    构建阶段基础镜像 & JDK 21 \\
    运行阶段基础镜像 & JRE 21 \\
    运行用户 & 非 root \\
    测试失败处理 & 阻止构建与部署 \\
    \hline
  \end{tabular}
\end{table}
